% ============================================================================
%  The Knowability--Capacity threshold for sign-of-benefit
%  (graded refinement of the binary impossibility, Thm 1 / Thm 2).
%  1-D Gaussian-location covariate-shift two-point model.
%
%  Status of claims (read this first):
%   * Lemma (benefit identity)               : PROVEN (exact, elementary).
%   * Theorem (population capacity, tau=1)    : PROVEN. Converse AND achievability
%       meet EXACTLY at the single threshold tau=1 for the population/identifiability
%       problem (Q_X known). This is a true capacity for that problem.
%   * Proposition (finite-n Le Cam boundary)  : PROVEN (finite-n lower bound). It
%       adds an O(1/sqrt n) layer around the threshold LOCATION, not a gap in the
%       threshold VALUE.
%   * Honest scope: the capacity is for THIS model and for sign-of-benefit; it is
%       NOT a general-distribution capacity, and tau=1 is an artifact of measuring
%       the drift budget eps in target-mass units (which is the natural,
%       label-free-meaningful choice). See Remark (scope).
%
%  All numbers validated by experiments/kbound/theory_validation/
%  val_knowability_capacity.py  ->  results_knowability_capacity.json (fixed seed).
% ============================================================================
\section{A knowability--capacity threshold for sign-of-benefit}\label{sec:knowcap}

Theorem~\ref{thm:imp} is \emph{binary}: it certifies a regime in which no
label-free rule beats chance. This section sharpens that dichotomy into a
\emph{graded, computable invariant} $K$ with an explicit threshold $\tau$, in a
single fully specified model, and proves that for the population (identifiability)
problem the converse and the achievability meet \emph{exactly} at $\tau$. Thus, in
this model, $K$ is a genuine \emph{capacity} for label-free sign recovery, not
merely a sufficient or a necessary condition. A finite-sample Le Cam analysis then
locates the threshold up to an $O(1/\sqrt n)$ boundary layer.

\subsection{Model (the concept mechanism is explicit)}\label{sec:knowcap-model}
Binary $Y\in\{0,1\}$, $0/1$ loss. Source covariates $P_X=\mathcal N(0,1)$ (known);
target covariates $Q_X=\mathcal N(\mu,1)$ (covariate shift; $\mu$ unknown, to be
estimated from an unlabeled target sample). The deployed classifier is
$g=f_0=\mathbf 1[x>0]$ (boundary at the origin). We are handed a \emph{fixed}
candidate adapted rule $f_a=\mathbf 1[x>a]$ with $a>0$ known --- the shipped
``adapted boundary'' --- and ask whether switching $g\to f_a$ helps. The
disagreement region is the interval $D=(0,a)$.

The label/concept mechanism must be explicit, because $\sign\Delta$ depends on
$P(Y\mid X)$, which $Q_X$ alone does not identify. We take \emph{hard-threshold
concepts} $\eta_\theta(x)=P(Y=1\mid x)=\mathbf 1[x>\theta]$. The
source-calibrated threshold $\theta_S$ is known (it is what a source-fit model
reports). The \emph{admissible target concepts} are those whose target
\emph{boundary mass} drifts from the source boundary by at most a budget
$\varepsilon\ge0$:
\begin{equation}\label{eq:knowcap-class}
  \mathcal C_\varepsilon \;=\;
  \Big\{\,\theta_T:\ \big|\,\Phi(\theta_T-\mu)-\Phi(\theta_S-\mu)\,\big|\le\varepsilon\,\Big\}.
\end{equation}
Here $\varepsilon$ is the $1$-D analogue of the calibration-drift budget $\beta$ of
the benefit-sign frontier (\S\ref{sec:bsfrontier}); by
Theorem~\ref{thm:conj1-dichotomy} some such untestable supplement is unavoidable,
and \eqref{eq:knowcap-class} fixes it in the natural label-free-meaningful unit of
\emph{target probability mass}. The label-free observer sees only
$(P_X,g,f_a,\theta_S,\varepsilon)$ and the unlabeled sample $\{x_i\}\sim Q_X$, and
must output a guess of $\sign\Delta$, where
\(
  \Delta=R_Q(f_0)-R_Q(f_a)
\)
($\Delta>0\iff$ adapting to $f_a$ helps).

\begin{lemma}[Exact benefit identity and the sign-flip locus]\label{lem:knowcap-identity}
On $D=(0,a)$ exactly one of $f_0,f_a$ is correct, so for the concept $\eta_\theta$,
\[
  \Delta(\theta)
  \;=\; Q\!\big((0,a)\cap\{x<\theta\}\big)-Q\!\big((0,a)\cap\{x>\theta\}\big)
  \;=\; 2\,\Phi(\theta'-\mu)-\Phi(-\mu)-\Phi(a-\mu),
\]
with $\theta'=\operatorname{clip}(\theta,0,a)$. $\Delta(\cdot)$ is nondecreasing in
$\theta$ with a single root
\[
  \boxed{\ \theta_0\;=\;\mu+\Phi^{-1}\!\Big(\tfrac12\big[\Phi(-\mu)+\Phi(a-\mu)\big]\Big)\ }
  \qquad(\text{the $Q$-median of the band }D),
\]
so $\sign\Delta=\sign(\theta-\theta_0)$. Both $\theta_0$ and the plug-in sign
$\sign(\theta_S-\theta_0)$ are \emph{observable}.
\end{lemma}
\begin{proof}
Off $D$, $f_0=f_a$, so the pointwise excess loss vanishes. On $D=(0,a)$ we have
$f_0=1,f_a=0$; for $0/1$ loss the excess loss of $f_0$ over $f_a$ is
$\mathbf 1[f_a{=}Y]-\mathbf 1[f_0{=}Y]=\mathbf 1[Y{=}0]-\mathbf 1[Y{=}1]=1-2Y$, with
conditional mean $1-2\eta_\theta(x)$, which is $+1$ for $x<\theta$ and $-1$ for
$x>\theta$. Integrating against $Q$ over $(0,a)$ gives the displayed difference of
band masses, hence
$\Delta(\theta)=\Phi(\theta'\!-\mu)-\Phi(-\mu)-\big(\Phi(a-\mu)-\Phi(\theta'\!-\mu)\big)
=2\Phi(\theta'\!-\mu)-\Phi(-\mu)-\Phi(a-\mu)$. This is nondecreasing in $\theta'$
(hence in $\theta$); setting it to zero gives $\Phi(\theta_0-\mu)=\tfrac12[\Phi(-\mu)+\Phi(a-\mu)]$,
the $Q$-median of $D$. (Validated to $4\times10^{-4}$ Monte-Carlo error against the
empirical $R_Q(f_0)-R_Q(f_a)$, \texttt{val\_knowability\_capacity.py} part (D).)
\end{proof}

\subsection{The computable invariant and the threshold}\label{sec:knowcap-K}
\begin{definition}[Knowability invariant]\label{def:knowcap-K}
\[
  K \;:=\; \frac{\big|\,\Phi(\theta_0-\mu)-\Phi(\theta_S-\mu)\,\big|}{\varepsilon}
       \;=\; \frac{\big|\,\tfrac12[\Phi(-\mu)+\Phi(a-\mu)]-\Phi(\theta_S-\mu)\,\big|}{\varepsilon}.
\]
$K$ is the \emph{mass margin} between the calibrated boundary $\theta_S$ and the
sign-flip locus $\theta_0$, measured in target probability mass, normalised by the
drift budget $\varepsilon$. It is a margin-weighted SNR of the covariate shift
relative to $g$'s boundary, computable from $(\mu,a,\theta_S,\varepsilon)$ alone
(with $\mu$ estimable label-free from $\{x_i\}$).
\end{definition}

\begin{theorem}[Population knowability--capacity, $\tau=1$]\label{thm:knowcap}
Assume $Q_X$ (equivalently $\mu$) is known. With threshold $\tau=1$:
\begin{enumerate}\itemsep1pt
\item[\textnormal{(i)}] \textbf{Achievability.} If $K>1$ then for \emph{every}
admissible target concept $\theta_T\in\mathcal C_\varepsilon$,
$\sign\Delta=\sign(\theta_S-\theta_0)$; the plug-in certificate
$\hat g=\sign(\theta_S-\theta_0)$ recovers the sign with error $0$.
\item[\textnormal{(ii)}] \textbf{Converse.} If $K<1$ there exist
$\theta_T^+,\theta_T^-\in\mathcal C_\varepsilon$ with \emph{identical} $Q_X$ (hence
$\TV=0$ between the laws of \emph{any} statistic of any number $n$ of unlabeled
observations) and $\sign\Delta(\theta_T^+)=-\sign\Delta(\theta_T^-)\neq0$.
Consequently the minimax label-free error equals $\tfrac12$ for every $n$.
\end{enumerate}
The converse and the achievability therefore meet at the \emph{single} threshold
$\tau=1$, \emph{exactly} (not asymptotically): $\{K>1\}$ is precisely the
identifiable set.
\end{theorem}
\begin{proof}
By Lemma~\ref{lem:knowcap-identity}, $\sign\Delta(\theta_T)=\sign(\theta_T-\theta_0)$
with $\theta_0$ fixed by the observables. The map $\theta\mapsto\Phi(\theta-\mu)$ is
a strictly increasing bijection, so the admissible set \eqref{eq:knowcap-class} is
exactly the interval
$\theta_T\in[\,\mu+\Phi^{-1}(p_S-\varepsilon),\,\mu+\Phi^{-1}(p_S+\varepsilon)\,]$ with
$p_S=\Phi(\theta_S-\mu)$, i.e.\ in mass coordinates $u=\Phi(\theta_T-\mu)$ it is the
interval $[\,p_S-\varepsilon,\,p_S+\varepsilon\,]$ (clipped to $[0,1]$). The sign of
$\Delta$ flips exactly at the mass level $u_0=\Phi(\theta_0-\mu)=\tfrac12[\Phi(-\mu)+\Phi(a-\mu)]$.

\emph{(i)} If $K>1$ then $|u_0-p_S|>\varepsilon$, so $u_0$ lies strictly outside
$[p_S-\varepsilon,p_S+\varepsilon]$. Hence every admissible $u\in[p_S-\varepsilon,p_S+\varepsilon]$
is on one side of $u_0$, so $\theta_T-\theta_0$ has constant sign over
$\mathcal C_\varepsilon$, equal to $\sign(\theta_S-\theta_0)=\sign(p_S-u_0)$. The
plug-in certificate is correct for all admissible concepts.

\emph{(ii)} If $K<1$ then $|u_0-p_S|<\varepsilon$, so $u_0\in(p_S-\varepsilon,p_S+\varepsilon)$
lies strictly interior to the admissible mass interval. Pick
$u^\pm=\Phi(\theta_T^\pm-\mu)$ with $u^+\in(u_0,p_S+\varepsilon]$ and
$u^-\in[p_S-\varepsilon,u_0)$ (both nonempty and admissible); then
$\sign\Delta(\theta_T^+)=+1$, $\sign\Delta(\theta_T^-)=-1$. Both worlds share the
\emph{same} $Q_X=\mathcal N(\mu,1)$ (the concept lives in the unobserved
$P(Y\mid X)$ channel), so the law of any unlabeled statistic is identical across
the two worlds: $\TV=0$ for every $n$. By Le Cam's two-point inequality the minimax
error of any label-free rule is $\ge\tfrac12(1-\TV)=\tfrac12$, and the constant
rules achieve $\tfrac12$, so it equals $\tfrac12$.

Finally, $K=1$ is the single boundary $|u_0-p_S|=\varepsilon$, where $u_0$ is an
endpoint of the admissible interval; the sign is constant on the (closed) interval
but not on its interior-plus-the-flip, so the transition is exactly at $\tau=1$.
(Validated: over $40{,}000$ random draws of $(\mu,a,\theta_S,\varepsilon)$, the
equivalence ``$K>1\iff$ identifiable over $\mathcal C_\varepsilon$'' holds with
\emph{zero} mismatches off a $10^{-6}$ shell; a boundary stress-test confirms the
flip at $K=1$ to $10^{-3}$. \texttt{val\_knowability\_capacity.py} part (A).)
\end{proof}

\smallskip\noindent\emph{Relation to the frontier.}
Theorem~\ref{thm:knowcap} is the \emph{graded} reading of
Theorem~\ref{thm:frontier}: the frontier's qualitative ``$|M|>\beta$'' becomes a
\emph{scalar} $K=|M|/\beta$-type ratio with the explicit value $\tau=1$ in a fully
specified location model, and the $0$-vs-$\tfrac12$ minimax dichotomy is recovered
as the two faces $K>1$ and $K<1$. The novelty here is that $K$ is shown to be
\emph{both} necessary and sufficient at the \emph{same} threshold --- a capacity ---
for the population problem, with $\theta_0$ (the flip locus) given in closed form by
the covariate shift $\mu$ and the candidate boundary $a$.

\subsection{Finite-$n$: a Le Cam boundary layer around the threshold}\label{sec:knowcap-finiten}
With only $n$ unlabeled samples, $\mu$ --- hence $K$ --- must be estimated, and the
deployable rule is the finite-$n$ \emph{three-way certificate}: estimate
$\hat\mu=\bar x$, form a $z$-confidence interval $[\hat\mu\pm z/\sqrt n]$, and commit
to $\sign(\theta_S-\hat\theta_0)$ only if $K$ stays above $1$ across the whole
interval (a lower-confidence bound $\hat K_{\mathrm{lcb}}>1$); otherwise abstain.
The next proposition shows the price of estimating $\mu$ is exactly an
$O(1/\sqrt n)$ layer around the threshold \emph{location}, with the same Le
Cam/Bretagnolle--Huber form as Proposition~\ref{prop:lecam-finite}.

\begin{proposition}[Finite-$n$ Le Cam boundary layer]\label{prop:knowcap-finiten}
Let $\mu^\star$ be the flip locus, $\theta_0(\mu^\star)=\theta_S$ (so $K=0$ there),
and fix the concept to transfer ($\theta_T=\theta_S$, i.e.\ $\varepsilon$ irrelevant;
the only noise is in estimating $\mu$). For $c>0$ consider the two worlds
$\mu^\pm=\mu^\star\pm\tfrac{c}{2}\sqrt{1/n}$, which straddle the flip and hence have
\emph{opposite} $\sign\Delta$. The unlabeled laws satisfy
\[
  \TV\big(\mathcal N(\mu^-,1)^{\otimes n},\mathcal N(\mu^+,1)^{\otimes n}\big)=2\Phi(c/2)-1,
  \qquad
  \mathrm{KL}=\tfrac{n}{2}\,(\mu^+-\mu^-)^2=\tfrac{c^2}{2},
\]
so every label-free sign test has minimax error
\[
  \inf_{\hat g}\max_{\pm}\Pr\!\big[\hat g\ \text{wrong}\big]
  \;\ge\;\tfrac12\big(1-\TV\big)=\Phi(-c/2)
  \;\ge\;\tfrac14\,e^{-c^2/2}\quad(\text{Bretagnolle--Huber}),
\]
\emph{constant in $n$}. Hence sign recovery is impossible below error $\Phi(-c/2)$
inside a window of half-width $\tfrac{c}{2}\sqrt{1/n}=\Theta(1/\sqrt n)$ around
$\mu^\star$, \emph{even when the concept transfers}.
\end{proposition}
\begin{proof}
The sample mean $\bar X\sim\mathcal N(\mu^\pm,1/n)$ is sufficient; the TV of two
equal-variance Gaussians with mean gap $\mu^+-\mu^-=c/\sqrt n$ and standard
deviation $1/\sqrt n$ is $2\Phi\big(\tfrac{c/\sqrt n}{2/\sqrt n}\big)-1=2\Phi(c/2)-1$,
and the per-sample $\mathrm{KL}\big(\mathcal N(\mu^+,1)\|\mathcal N(\mu^-,1)\big)=\tfrac12(\mu^+-\mu^-)^2$
tensorises to $\tfrac{c^2}{2}$. A sign test is a test between the two worlds; Le
Cam's testing affinity gives the $\tfrac12(1-\TV)$ floor and Bretagnolle--Huber
gives the $\tfrac14 e^{-\mathrm{KL}}$ certificate. That the two worlds straddle the
flip (so the signs are genuinely opposite) is Lemma~\ref{lem:knowcap-identity}.
(Validated: the empirical Bayes-optimal sign-test error meets $\Phi(-c/2)$ to
$\le1.5\times10^{-2}$, exceeds $\tfrac14 e^{-c^2/2}$, and is constant across
$n\in\{100,1000,10000\}$ for $c\in\{0.5,1,2\}$;
\texttt{val\_knowability\_capacity.py} part (C).)
\end{proof}

\noindent The finite-$n$ certificate \emph{matches} this: numerically its minimax
sign-recovery accuracy $\to1$ wherever $K_{\mathrm{pop}}>1$ and $=\tfrac12$ wherever
$K_{\mathrm{pop}}<1$, with the transition pinned at $K=1$, across
$n\in\{200,2000,20000\}$ (Table~\ref{tab:knowcap-phase};
\texttt{val\_knowability\_capacity.py} part (B)). The boundary layer is therefore an
$O(1/\sqrt n)$ blur of the threshold \emph{location} $\mu^\star$, not a gap in the
threshold \emph{value} $\tau=1$.

\begin{table}[t]\centering\small
\begin{tabular}{rrccc}
\hline
$\mu$ & $K_{\mathrm{pop}}$ & acc $n{=}200$ & acc $n{=}2000$ & acc $n{=}20000$\\
\hline
$-0.807$ & $1.416$ & $0.994$ & $1.000$ & $1.000$\\
$\ \ 0.366$ & $1.679$ & $0.995$ & $1.000$ & $1.000$\\
$\ \ 1.633$ & $1.677$ & $0.995$ & $1.000$ & $1.000$\\
$\ \ 2.806$ & $1.417$ & $0.994$ & $1.000$ & $1.000$\\[2pt]
$-1.407$ & $0.640$ & $0.500$ & $0.500$ & $0.500$\\
$\ \ 0.980$ & $0.063$ & $0.500$ & $0.500$ & $0.500$\\
$\ \ 1.000$ & $0.000$ & $0.500$ & $0.500$ & $0.500$\\
$\ \ 3.406$ & $0.641$ & $0.500$ & $0.500$ & $0.500$\\
\hline
\end{tabular}
\caption{Phase transition of the finite-$n$ three-way certificate
($a{=}2,\theta_S{=}1,\varepsilon{=}0.05$; minimax over the admissible concept). Above
$\tau{=}1$ the minimax accuracy $\to1$; below it the certificate abstains and the
accuracy is exactly $\tfrac12$. (Real numbers, seed fixed;
\texttt{results\_knowability\_capacity.json}.)}
\label{tab:knowcap-phase}
\end{table}

\begin{remark}[Honest scope --- what is and is not claimed]\label{rmk:knowcap-scope}
\textnormal{(1)} Theorem~\ref{thm:knowcap} is a capacity \emph{for this $1$-D
Gaussian-location model and for sign-of-benefit only}; it is not a
general-distribution capacity. \textnormal{(2)} The clean value $\tau=1$ is a
consequence of measuring the drift budget $\varepsilon$ in target-\emph{mass} units
\eqref{eq:knowcap-class}; this is the natural label-free-meaningful unit (it makes
$\theta_0$ the $Q$-median and the flip condition a mass comparison), but a different
budget parametrisation would rescale $\tau$. \textnormal{(3)} The same equivalence
$K>1\iff$ identifiable holds for \emph{soft} (logistic) concepts of fixed known
slope, with $\theta_0$ the (numerically computed) root of $\Delta(\cdot)$ rather than
the band median --- verified with zero mismatches over $1{,}496$ draws --- so the
result is not an artifact of the hard-threshold concept; only the closed form for
$\theta_0$ is special to it. \textnormal{(4)} Relative to the existing five
theorems, this is a \emph{refinement}, not a new pillar: it converts the
$0$-vs-$\tfrac12$ frontier of Theorem~\ref{thm:frontier} into a scalar capacity
with an explicit closed-form flip locus, and supplies the finite-$n$ boundary layer
that Proposition~\ref{prop:lecam-finite} left implicit for the location family. Its
value is a sharpened, fully worked \emph{instance}, useful pedagogically and as the
tight test-bed for the graded invariant; it does not enlarge the set of
\emph{distribution-free} guarantees.
\end{remark}
